\begin{summarybox}
\textbf{一句话理解：} 高斯函数分成两部分看：$\lfloor x\rfloor$ 是整数部分，$\{x\}=x-\lfloor x\rfloor$ 是小数部分。
\[\boxed{\lfloor x\rfloor=n\iff n\le x<n+1\quad(n\in\mathbb Z).}\]
\[\boxed{\{x\}=x-\lfloor x\rfloor,\quad 0\le \{x\}<1.}\]
\[\begin{array}{c|c}\text{取整条件}&\text{普通不等式}\\\hline \lfloor x\rfloor\ge n&x\ge n\\\lfloor x\rfloor<n&x<n\\\lfloor x\rfloor\le n&x<n+1\\\lfloor x\rfloor>n&x\ge n+1\end{array}\]
\textbf{图像对照：} 两者定义域均为 $\mathbb R$；整数部分值域为 $\mathbb Z$，呈台阶型；小数部分值域为 $[0,1)$，呈周期性锯齿型。
\textbf{四条证明线索：} 最大整数定义、整数平移与小数部分、整数边界转换、整数与小数部分分解。
\[\lfloor x+m\rfloor=\lfloor x\rfloor+m\quad(m\in\mathbb Z).\]
\[\lfloor x\rfloor+\lfloor y\rfloor\le\lfloor x+y\rfloor\le\lfloor x\rfloor+\lfloor y\rfloor+1.\]
\textbf{考前提醒：} 负数向左取；区间左闭右开；加法可能进位；小数部分最小正周期为 $1$，取整函数本身没有正周期。
\end{summarybox}
